\documentclass[11pt]{article}
\usepackage{mathblatt}
% Präambel des Setzers (werkzeuge/setzer-praeambel.tex, wortgleich); dazu \gteil (Muster T2U)
\makeatletter
\setlength{\mbpfbe}{0mm}% keine Punkte (pruefung.md § 5: nur beim Prüfstein)
\newlength{\gEin}\setlength{\gEin}{0pt}
\let\gSchrift\relax
\renewcommand{\pfkreuzzeile}[1]{\par\noindent\hangindent3.2em\hangafter1\hspace*{1.6em}$\square$\ #1\par}
\newcommand{\gkreuz}[1]{$\square$\ #1\hspace{2em}}
\newcommand{\gkreuzzeile}[2]{\par\noindent\hangindent2.2em\hangafter1\mbox{#1\,$\square$}\ #2\par}% Buchstabe, Option: eine je Zeile (neben Skizze)
% links, was man ansieht, rechts, was man tut; Spaltengrenze überall an derselben Stelle (gemeinsam.md „Satz“)
\newsavebox{\gbox}\newlength{\gLinks}
% v0.10 (Lehrer 13.10.): keine starre Halbierung – die linke Spalte ist so breit wie die Figur (+ 6 mm), die
% Antwort steht direkt daneben; ist die Figur breiter als 0,62 Zeile, steht die Antwort darunter
\newcommand{\gzwei}[2]{\par\smallskip\noindent\sbox{\gbox}{#1}%
  \setlength{\gLinks}{\dimexpr\wd\gbox+6mm\relax}%
  \ifdim\gLinks>0.62\linewidth
    \usebox{\gbox}\par\smallskip\noindent\begin{minipage}{\linewidth}#2\end{minipage}%
  \else
    \begin{minipage}[t]{\gLinks}\vspace{0pt}\usebox{\gbox}\end{minipage}%
    \begin{minipage}[t]{\dimexpr\linewidth-\gLinks\relax}\vspace{0pt}\raggedright #2\end{minipage}%
  \fi\par}
% v0.9 (12.10., pruefung.md § 5 „Marke mit B, Bedingungen rechts“): drittes Argument von pfaufg = Bedingungen
% der Prüfung (v0.10: \pfohnetr = TR durchgestrichen, „FOR“), klein grau im rechten Rand auf der Grundlinie der Nummer, Zeilen mit \\
\newcommand{\gbedrand}[2]{\ifblank{#2}{}{\rlap{\hspace*{\dimexpr#1+2mm\relax}%
  \parbox[t]{15mm}{\raggedright\scriptsize\color{mbgrau}\strut#2}}}}
% Teil eines Bündels (a, b …): Marke links im Rand, Bedingungen rechts im Rand, auf der Zeile des Teils
\newcommand{\gteilrand}[2]{\llap{\mbpfrandmarke{#1}\hspace*{\mbpfnr}}\gbedrand{\linewidth}{#2}}
% Teil eines Bündels: Buchstabe in eigener Spalte, Inhalt eingerückt (\gEin = Einzug, damit die Spaltengrenze von
% \gzwei an derselben Stelle bleibt)
\newenvironment{gbteil}[3]{\par\smallskip\noindent\gteilrand{#2}{#3}\makebox[\mbpfteil][l]{\textbf{#1}}%
  \begin{minipage}[t]{\dimexpr\linewidth-\mbpfteil\relax}\setlength{\gEin}{\mbpfteil}\raggedright\strut\ignorespaces}%
 {\end{minipage}\par}
% Grafik höchstens so breit wie die Spalte (Bündel nebeneinander, Koordinatensystem ohne Karo)
\newsavebox{\gpbox}
\newcommand{\gpassend}[1]{\sbox{\gpbox}{#1}\ifdim\wd\gpbox>\linewidth\typeout{GPASSEND: verkleinert \the\wd\gpbox\space auf \the\linewidth}\resizebox{\linewidth}{!}{\usebox{\gpbox}}\else\usebox{\gpbox}\fi}
% Skizzenfeld (v0.10, Lehrer 13.10.): Karo ab dem linken Rand über die Zeile, „Skizze“ klein grau oben links im Feld
% (wie \pfskizzenfeld der Vorlage); der rechte Rand bleibt den Bedingungen
\newcommand{\gskizze}[1]{\par\smallskip\noindent\begin{tikzpicture}
  \clip (0,0) rectangle ({\linewidth-1mm},{#1*10mm});
  \draw[black!16,line width=0.3pt,step=5mm] (0,0) grid ({\linewidth},{#1*10mm});
  \node[anchor=north west,inner sep=2pt,font=\scriptsize,text=mbgrau,fill=white] at (0,{#1*10mm}) {Skizze};
  \end{tikzpicture}\par}
\RenewDocumentEnvironment{pfaufg}{m m m}{%
  \begin{lrbox}{\mbaufgabebox}\begin{minipage}[t]{\linewidth}%
  \gSchrift\noindent\hspace*{\gEin}\mbpfrandmarke{#1}\gbedrand{\linewidth-\gEin}{#3}%
  \begin{minipage}[t]{\mbpfnr}\strut\textbf{#2}\end{minipage}%
  \begin{minipage}[t]{\dimexpr\linewidth-\gEin-\mbpfnr\relax}\raggedright\strut\ignorespaces}%
 {\end{minipage}%
  \end{minipage}\end{lrbox}%
  \setlength{\mbaufgabehoehe}{\dimexpr\ht\mbaufgabebox+\dp\mbaufgabebox\relax}%
  \ifdefempty{\gkopfnach}{\par\addvspace{5pt}\Needspace*{\mbaufgabehoehe}}{%
    % Stufenkopf, der mit dieser (ersten) Aufgabe zusammenbleibt: Platz für beide verlangen, dann den Kopf setzen
    % (12.10., Quadratische: ein Kopf stand allein unten, die Aufgabe mit Koordinatensystem auf der nächsten Seite)
    \par\Needspace*{\dimexpr\mbaufgabehoehe+40pt\relax}\gkopfnach\gdef\gkopfnach{}\par\addvspace{5pt}}%
  \noindent\usebox{\mbaufgabebox}\par\addvspace{5pt}}
% Stufenkopf vormerken: \gkopfvor{<Kopf>} – die nächste pfaufg setzt ihn, sobald beide zusammen Platz haben
\def\gkopfnach{}
\newcommand{\gkopfvor}[1]{\gdef\gkopfnach{#1}}
% Teilaufgabe mit Skizze links in gzwei: Buchstabe in eigener Spalte, darunter Text und Skizze
\newcommand{\gteil}[2]{\makebox[\mbpfteil][l]{\textbf{#1}}\parbox[t]{\dimexpr0.5\textwidth-\gEin-\mbpfnr-5mm-\mbpfteil\relax}{\raggedright #2}}
% Stufenkopf der Vorlage (mathblatt.sty 2026-10-12: Linie, Abstand, 13 pt, \pfjahre) ohne ihren Vorhalt von
% 14 Zeilen: der Setzer fordert vor dem Kopf 0,16 Seitenhöhe (Kopf + erste Aufgabe), mehr nicht
\setlength{\mbpfstufevorhalt}{0pt}
\makeatother
\newcommand{\karo}[1]{\par\smallskip\noindent\begin{tikzpicture}
  \clip (0,0) rectangle ({\linewidth-1mm},{#1*10mm});
  \draw[black!16,line width=0.3pt,step=5mm] (0,0) grid ({\linewidth},{#1*10mm});
  \end{tikzpicture}\par}
\newcommand{\kk}[1]{$\square$\,#1\hspace{0.9em}}
\newcommand{\linie}[1]{\,{\color{black!45}\rule[-1pt]{#1}{0.4pt}}\,}
\newcommand{\pkt}{\textperiodcentered{}}
% v0.7 (12.10.): Schülerfrage unter dem Stufenkopf; Ankreuzoption als Bild (Buchstabe und Kästchen darunter)
\newcommand{\gfrage}[1]{\par\nobreak\noindent\hspace*{\mbpfnr}{\small\itshape #1}\par\nobreak\addvspace{2pt}}
\newlength{\gbildbreite}\setlength{\gbildbreite}{0.32\linewidth}% Setzer setzt sie je Reihe (vier Bilder: 0,24)
\newcommand{\gbildoption}[3]{\begin{minipage}[t][#3][b]{\gbildbreite}\centering #2\par\smallskip\mbox{#1\,$\square$}\end{minipage}}
% Lineare (12.10.): Bild mit Buchstabe ohne Kästchen (Diagramm zuordnen); Aussage mit „wahr / falsch“ zum Ankreuzen
\newcommand{\gbildname}[3]{\begin{minipage}[t][#3][b]{\gbildbreite}\centering #2\par\smallskip\mbox{\textbf{#1}}\end{minipage}}
\newcommand{\gwahrfalsch}[2]{\par\noindent\hangindent2.2em\hangafter1\makebox[2.2em][l]{(#1)}#2\hfill\mbox{wahr $\square$\quad falsch $\square$}\par}
\begin{document}

\pfheftstil
\pflaufkopf[P10 \pkt{} Prüfung]{Gleichungssystem im Sachzusammenhang}{93F1}\pfseitevon
\pfheftkopf{Gleichungssystem im Sachzusammenhang}{}
\pfstufenzeile{Gleichungen aufstellen \pkt{} x und y gesucht \pkt{} Bedeutung gesucht}

% ===================== Stufe: aufstellen =====================
\gkopfvor{\pfstufekopf{Gleichungen aufstellen}{\pfjahre{24}{21 \pkt{} 16}}}
% kein Einstieg: keine Schülerfrage in der Gliederung (Schülerfrage, -Fälle, -Form, -Weg, -Skizze; oder Frage) – Einstieg nur von dort, nie aus der Bank (pruefung.md § 5, 12.10.)
% 2024-OS-K7a-m1 (Original 2024-OS-K7a, Schwere I, rang 1; herausgelöste Fassung: keine)
\begin{pfaufg}{nach~P10\\’24~\pkt{}~7a}{1.}{}
Die Brüder Daniel und Filip kaufen ihrer Mutter zum Muttertag je einen Blumenstrauß, Daniel in Laden A, Filip in Laden B. In Laden A zahlt man für eine Rose und eine Tulpe zusammen 3,60 €. Daniel bezahlt für 7 Rosen und 4 Tulpen insgesamt 20,70 €. Stelle ein lineares Gleichungssystem zu diesen Angaben auf ($x$: Preis einer Rose, $y$: Preis einer Tulpe).
\par\smallskip\par\noindent \hfill I.: \linie{30mm}\hphantom{\,cm²}\par\noindent \hfill II.: \linie{30mm}\hphantom{\,cm²}
\fusshilfe{\mbox{1:~3,60; 20,70}}
\end{pfaufg}

% ===================== Stufe: lösen =====================
% Formel-Frage (einmal je Blatt, erste Stufe mit Formel; pruefung.md § 5)
\Needspace*{0.16\textheight}\pfstufekopf{x und y gesucht}{\pfjahre{22}{}}\gfrage{Schlag zuerst die Formel nach. Formelsammlung, Seite~\linie{12mm}}
% Schülerfrage der Stufe (Gliederung; Einstieg nur von dort, nie aus der Bank – 12.10.)
\begin{pfaufg}{}{2.}{}
Beim Einsetzen stellst du eine Gleichung um und setzt sie in die andere ein. Fang so an:
\par\smallskip
\par\noindent\makebox[\mbpfteil][l]{\textbf{a)}}Stelle x + y = 10 nach y um\hfill\linie{30mm}
\par\noindent\makebox[\mbpfteil][l]{\textbf{b)}}Setze y = 10 $-$ x in 3x + 2y = 24 ein\hfill\linie{30mm}
\par\fusshilfe{\mbox{2:~a) y = 10 $-$ x; b) 3x + 2 · (10 $-$ x) = 24}}
\end{pfaufg}
% 2022-OS-K7b-m1 (Original 2022-OS-K7b, Schwere II, rang 3; herausgelöste Fassung: keine)
\begin{pfaufg}{nach~P10\\’22~\pkt{}~7b}{3.}{FOR}
Frau Bauer will in ihrem Garten Obstbäume pflanzen. Zusammen kosten ein Apfelbaum und ein Pflaumenbaum 58,00 €. Sie kauft 8 Apfelbäume und 4 Pflaumenbäume und zahlt dafür 330,40 €. Sie schreibt dazu zwei Gleichungen auf: I $x + y = 58$; II $8x + 4y = 330{,}40$. Berechne, was ein Apfelbaum und was ein Pflaumenbaum kostet.
\par\smallskip\par\noindent \hfill Apfelbaum: \linie{30mm}\,€\par\noindent \hfill Pflaumenbaum: \linie{30mm}\,€
\fusshilfe{\mbox{3:~Apfelbaum 24,60 €; Pflaumenbaum 33,40 €}}
\end{pfaufg}

% ===================== Stufe: deuten =====================
\gkopfvor{\pfstufekopf{Bedeutung gesucht}{\pfjahre{24 \pkt{} 22}{}}}
% kein Einstieg: keine Schülerfrage in der Gliederung (Schülerfrage, -Fälle, -Form, -Weg, -Skizze; oder Frage) – Einstieg nur von dort, nie aus der Bank (pruefung.md § 5, 12.10.)
% 2022-OS-K7a-m1 (Original 2022-OS-K7a, Schwere I, rang 2; herausgelöste Fassung: keine)
\begin{pfaufg}{nach~P10\\’22~\pkt{}~7a}{4.}{}
Frau Bauer will in ihrem Garten Obstbäume pflanzen. Zusammen kosten ein Apfelbaum und ein Pflaumenbaum 58,00 €. Sie kauft 8 Apfelbäume und 4 Pflaumenbäume und zahlt dafür 330,40 €. Sie schreibt dazu zwei Gleichungen auf: I $x + y = 58$; II $8x + 4y = 330{,}40$. Gib an, wofür die Variablen $x$ und $y$ in Frau Bauers Gleichungen stehen.
\par\smallskip\par\noindent \hfill x steht für: \linie{30mm}\hphantom{\,cm²}\par\noindent \hfill y steht für: \linie{30mm}\hphantom{\,cm²}
\fusshilfe{\mbox{4:~Preis eines Apfelbaums in €; Preis eines Pflaumenbaums in €}}
\end{pfaufg}
% 2024-OS-K7b-m1 (Original 2024-OS-K7b, Schwere II, rang 4; herausgelöste Fassung: keine)
\begin{pfaufg}{nach~P10\\’24~\pkt{}~7b}{5.}{FOR}
Die Brüder Daniel und Filip kaufen ihrer Mutter zum Muttertag je einen Blumenstrauß, Daniel in Laden A, Filip in Laden B. In Laden B kostet eine Rose 2,40 € und eine Tulpe 1,60 €. Filip schreibt zu seinem Strauß das Gleichungssystem I. $2{,}40r + 1{,}60t = 28{,}80$; II. $r + t = 15$. Schreib in einem Satz auf, was Gleichung II über den Strauß aussagt. Löse das Gleichungssystem.
\par\smallskip\par\noindent \hfill Gleichung II bedeutet: \linie{30mm}\hphantom{\,cm²}\par\noindent \hfill Rosen: \linie{30mm}\hphantom{\,cm²}\par\noindent \hfill Tulpen: \linie{30mm}\hphantom{\,cm²}
\fusshilfe{\mbox{5:~Filip kauft zusammen 15 Blumen; 6 Rosen; 9 Tulpen}}
\end{pfaufg}

\pfblattende{11}{26 & 25 & 24 & 23 & 22 & 21 & 20 & 19 & 18 & 17 & 16}{ &  & 7a 7b &  & 7a 7b & 7b &  &  &  &  & 6d}%
{\pfausschnitt{}{Gleichungssystem im Sachzusammenhang}{}}
\end{document}